% TOP_PROBLEM: {"id": 2800102, "title": "10 Lectures and 42 Open Problems — Gaussian singular-value monotonicity", "category_id": 19, "category": {"id": 19, "name": "probability", "display_name": "Probability", "description": "Problems involving probability theory, stochastic processes, and random structures.", "slug": "probability", "order_index": 19, "created_at": "2026-07-31T15:26:25.671Z"}, "status": "open"}
% FIRST_POSTED: null
% ENTRY_KIND: statement_only
% Imported from: batch08.tex; UnsolvedMath ID 2800102
% Compile twice: lualatex 2800102.tex
\documentclass[11pt]{article}
\usepackage{fontspec}
\setmainfont{Latin Modern Roman}
\usepackage{amsmath,amssymb,mathtools,unicode-math}
\setmathfont{Latin Modern Math}
\usepackage{xurl,hyperref}
\hypersetup{hidelinks,pdftitle={UnsolvedMath Problem 2800102}}
\newcommand{\sref}[2]{\hyperlink{source:#1:#2}{[S#2]}}
\newcommand{\cataloguescope}[1]{\paragraph{Mathematical question.}#1}
\begin{document}
% UnsolvedMath ID 2800102; source catalogue code AMR-027-0102
\setcounter{section}{2800101}
\section{10 Lectures and 42 Open Problems — Gaussian singular-value monotonicity}\label{2800102}
{\small\sffamily UnsolvedMath ID 2800102 · Probability}
\subsection{Definitions and mathematical statement}

\paragraph{Shared notation.}$\mathbb N=\{1,2,\ldots\}$, and $\mathbb Z,\mathbb Q,\mathbb R,\mathbb C$ denote the integers, rationals, reals and complex numbers. Any other conventions are specified locally.


For an integer \(d\ge1\), let \(G_{\mathbb R}^{(d)}\) have independent entries distributed as \(\mathcal N(0,1/d)\). Let \(G_{\mathbb C}^{(d)}\) have independent entries
\[\frac{A_{ij}+iB_{ij}}{\sqrt{2d}},\]
where all \(A_{ij},B_{ij}\) are independent \(\mathcal N(0,1)\) variables, so that each complex entry has mean zero and \(\mathbb E|G_{ij}|^2=1/d\). For either field \(\mathbb K\), write \(\sigma_1(G),\ldots,\sigma_d(G)\) for the nonnegative singular values, the square roots of the eigenvalues of \(G^*G\), and define
\[\alpha_{\mathbb K}(d)=\frac1d\mathbb E\sum_{j=1}^d\sigma_j\bigl(G_{\mathbb K}^{(d)}\bigr).\]
\paragraph{Statement recorded in the supplied source.}
Do both inequalities
\[\alpha_{\mathbb R}(d+1)\ge\alpha_{\mathbb R}(d),\qquad
  \alpha_{\mathbb C}(d+1)\le\alpha_{\mathbb C}(d)\]
hold for every integer \(d\ge1\)? \sref{2800102}{2}
\subsection{Short English statement}
Does the expected average singular value increase with matrix dimension for normalized real Gaussian matrices and decrease for normalized complex Gaussian matrices?
\subsection{Sources}
\begin{description}
\item[\hypertarget{source:2800102:1}{S1}] ULAM.AI and the UnsolvedMath Contributors, UnsolvedMath: A Curated Collection of Open Mathematics Problems, record 2800102 (AMR-027-0102). CC BY 4.0. \url{https://huggingface.co/datasets/ulamai/UnsolvedMath}.
\item[\hypertarget{source:2800102:2}{S2}] Afonso S. Bandeira, Christopher Kennedy and Amit Singer, Approximating the Little Grothendieck Problem over the Orthogonal and Unitary Groups (2013; revised 2016), arXiv:1308.5207, Definition 2 and Conjecture 8. \url{https://arxiv.org/pdf/1308.5207}.
\end{description}
\label{end:2800102}
\end{document}
